% Anexo — Fichas de registro de uso de artefactos y registro de usos de oficinas (AS1).
% Origen: 01_normativa/proyecto/originales/2026-01_ficha-registro-uso-artefactos-oficina.docx (ficha)
%         y 2026_calculo-produccion-aguas-servidas.docx (registro y cálculo), documentos del proyecto ANID.
% Modelo de redacción indicado por el profesor guía: tesis de Yévenes (2019), anexos 1–3 y 10–11
% (05_ejemplos/fichas/2019_ucm_yevenes_aguas-grises-rurales.md).
\chapter{Fichas de registro de uso de artefactos sanitarios}
\label{anx:fichas-registro}

Este anexo presenta las fichas utilizadas para registrar el uso de los artefactos sanitarios de las oficinas (AS1) y los registros obtenidos, que sustentan la estimación de la producción de aguas servidas descrita en la metodología (\cref{sec:met-caracterizacion}).
Las \cref{fig:ficha-lavaplatos,fig:ficha-lavamanos,fig:ficha-inodoro} muestran la ficha de cada artefacto, con una tabla por semana en la que los usuarios marcaron cada uso.
La \cref{tab:registro-usos-oficina} resume el número de usos diarios registrados.

\begin{figure}[htbp]
  \centering
  \includegraphics[page=1, width=0.5\textwidth, trim=1.5cm 2.5cm 1.5cm 1cm, clip]{fichas/ficha_registro_uso_oficina.pdf}
  \caption{Ficha de registro de uso del lavaplatos de las oficinas (semanas del 19 de enero al 1 de febrero de 2026). Fuente: proyecto ANID/Anillos/ATE250033.}
  \label{fig:ficha-lavaplatos}
\end{figure}

\begin{figure}[htbp]
  \centering
  \includegraphics[page=2, width=0.5\textwidth, trim=1.5cm 2.5cm 1.5cm 1cm, clip]{fichas/ficha_registro_uso_oficina.pdf}
  \caption{Ficha de registro de uso del lavamanos de las oficinas. Fuente: proyecto ANID/Anillos/ATE250033.}
  \label{fig:ficha-lavamanos}
\end{figure}

\begin{figure}[htbp]
  \centering
  \includegraphics[page=3, width=0.5\textwidth, trim=1.5cm 2.5cm 1.5cm 1cm, clip]{fichas/ficha_registro_uso_oficina.pdf}
  \caption{Ficha de registro de uso del inodoro de las oficinas (X: descarga grande; O: descarga pequeña). Fuente: proyecto ANID/Anillos/ATE250033.}
  \label{fig:ficha-inodoro}
\end{figure}

\begin{table}[htbp]
  \centering
  \caption{Número de usos diarios registrados por artefacto en las oficinas (AS1), de lunes a viernes. Fuente: proyecto ANID/Anillos/ATE250033.}
  \label{tab:registro-usos-oficina}
  \begin{tabular}{@{}llccccc@{}}
    \toprule
    & & \multicolumn{2}{c}{Enero (escenario~1)} & \multicolumn{3}{c}{Febrero (escenario~2)} \\
    \cmidrule(lr){3-4}\cmidrule(l){5-7}
    Artefacto & Día & Sem.~3 & Sem.~4 & Sem.~1 & Sem.~2 & Sem.~3 \\
    \midrule
    Inodoro & Lunes     & 8 & 8  & 6 & 4 & 5 \\
            & Martes    & 9 & 9  & 4 & 5 & 5 \\
            & Miércoles & 9 & 10 & 4 & 4 & 4 \\
            & Jueves    & 9 & 7  & 3 & 9 & 3 \\
            & Viernes   & 8 & 10 & 3 & 5 & --- \\
            & \textit{Promedio del escenario} & \multicolumn{2}{c}{\textit{9}} & \multicolumn{3}{c}{\textit{5}} \\
    \midrule
    Lavamanos & Lunes     & 7  & 10 & 6 & 4 & 5 \\
              & Martes    & 10 & 9  & 5 & 6 & 6 \\
              & Miércoles & 10 & 12 & 4 & 3 & 5 \\
              & Jueves    & 9  & 8  & 3 & 9 & 3 \\
              & Viernes   & 8  & 12 & 3 & 6 & --- \\
              & \textit{Promedio del escenario} & \multicolumn{2}{c}{\textit{10}} & \multicolumn{3}{c}{\textit{5}} \\
    \midrule
    Lavaplatos & Lunes     & 5 & 6   & 4 & 3 & 4 \\
               & Martes    & 5 & 5   & 3 & 3 & 5 \\
               & Miércoles & 6 & 4   & 6 & 4 & 4 \\
               & Jueves    & 4 & 8   & 4 & 5 & 5 \\
               & Viernes   & 7 & --- & 5 & 4 & 2 \\
               & \textit{Promedio del escenario} & \multicolumn{2}{c}{\textit{6}} & \multicolumn{3}{c}{\textit{4}} \\
    \bottomrule
    \multicolumn{7}{@{}p{0.9\textwidth}@{}}{\footnotesize ---: sin registro. Promedios redondeados al entero.}
  \end{tabular}
\end{table}
\revisar{En el documento original la cuarta columna aparece rotulada \enquote{02S1} (repetida); se interpretó como la segunda semana de febrero.
La ficha impresa cubre del 19 de enero al 1 de febrero, mientras que el registro llega a la tercera semana de febrero: confirmar si se usaron más fichas.}
